\section{Related Work：文献检索与精读}
\label{sec:related}

\webref{docs/2-文献/}
\resref{resources/02-literature/}

\subsection{检索渠道与来源优先级}
建议渠道包括：知网综述（英文适应前的地图）、Google Scholar、arXiv、AMiner、PubScholar、Papers with Code，并以 CCF 列表与组内投稿目标（见图~\ref{fig:journal-list}，亦见 Case Study）辅助判断权威性。优先浏览 PR、TPAMI、TIP、TNNLS、AAAI、CVPR、ECCV、ICCV、ACM MM、NeurIPS 等来源。

\subsection{筛选与文献 List}
策略上先读英文综述建立领域地图，再筛选经典/顶会顶刊/有代码有数据的工作；社区解读仅作辅助，结论必须回原文核对。每条文献记录应至少包含题目、作者、引用、时间、链接、代码与数据集字段。仓库提供 \path{resources/02-literature/literature_list.csv} 及示例清单。

\subsection{阅读顺序与汇报}
泛读顺序为 Abstract 与文中总览图、Introduction、Conclusion；确有价值再读 Method 与 Experiment。精读后应能陈述问题定义、背景不足、方法框架、实验结论与可迁移点。读懂第一篇领域论文常需一至两周，并连带阅读十余篇相关工作，属正常现象。论文汇报与周汇报 checklist 见网页文献篇；扫描要点图册见 \path{resources/02-literature/literature_and_report_tips.pdf}。文献管理推荐 Zotero。
